\section{The Light of Glory and the Vision of the Trinity}

Whoever weighs what the Church defines in \textit{Benedictus Deus} and \textit{Laetentur caeli} meets a towering mystery. Scripture says that God ``inhabiteth light inaccessible, whom no man hath seen, nor can see'' (1 Tim 6:16), and St.\ John writes, ``No man hath seen God at any time'' (John 1:18). How can a finite, created mind look directly upon the infinite Essence of the Father, the Son, and the Holy Ghost with no created image standing between them?

The Fathers felt the full weight of that question. Examining the Old Testament theophanies in Book III of \textit{On the Trinity}, St.\ Augustine insists that the burning bush, the pillar of cloud, and the appearances to Abraham occurred through created signs and angelic ministry, never by exposing the immutable Godhead to mortal sight: ``Wherefore the substance, or, if it is better so to say, the essence of God, wherein we understand, in proportion to our measure, in however small a degree, the Father, the Son, and the Holy Spirit, since it is in no way changeable, can in no way in its proper self be visible.''\endnote{St.\ Augustine, \textit{On the Trinity} III.11.21--22, trans.\ Arthur West Haddan (1887), revised by Kevin Knight.}

In the East, St.\ Gregory of Nazianzus---honored by the Church as ``the Theologian''---felt this transcendence as both awe and longing. In his Second Theological Oration, ascending Mount Sinai in spirit like Moses, he confesses the limit of mortal contemplation:

\begin{quote}
What is this that has happened to me, O friends, and initiates, and fellow-lovers of the truth? I was running to lay hold on God, and thus I went up into the Mount, and drew aside the curtain of the Cloud, and entered away from matter and material things, and as far as I could I withdrew within myself. And then when I looked up, I scarce saw the back parts of God; although I was sheltered by the Rock, the Word that was made flesh for us. And when I looked a little closer, I saw, not the First and unmingled Nature, known to Itself---to the Trinity, I mean; not That which abideth within the first veil, and is hidden by the Cherubim; but only that Nature, which at last even reaches to us.\endnote{St.\ Gregory of Nazianzus, \textit{Oration} 28.3 (Second Theological Oration), trans.\ Charles Gordon Browne and James Edward Swallow, NPNF2 vol.\ 7.}
\end{quote}

Yet Gregory insists that purity opens the soul ``in order that light may be apprehended by light,'' praying that ``one illumination may come upon us from the One God, One in diversity, diverse in Unity, wherein is a marvel'' (\textit{Oration} 28.1). In his Second Oration on Easter he ascends from vigil lamps and angelic splendor to the uncreated light ``in the Trinity, from which all light derives its being'' (\textit{Oration} 45.2). And in Oration 16 he states plainly what awaits the just: ``Some will be welcomed by the unspeakable light and the vision of the holy and royal Trinity, Which now shines upon them with greater brilliancy and purity and unites Itself wholly to the whole soul, in which solely and beyond all else I take it that the kingdom of heaven consists.''\endnote{St.\ Gregory of Nazianzus, \textit{Oration} 16.9, NPNF2 vol.\ 7.}

In the thirteenth century, St.\ Thomas Aquinas gave this patristic mystery its classical theological architecture in the \textit{Prima Pars} (\textit{ST} I, q.\ 12) and in his earlier commentary on the \textit{Sentences}, preserved in the posthumously compiled \textit{Supplement} (q.\ 92, a.\ 1). God is hidden from us not by obscurity, Aquinas begins, but by excess of intelligibility: being Pure Act, He is supremely knowable in Himself, yet dazzles finite eyes as the sun overwhelms the eyes of a bat (\textit{ST} I, q.\ 12, a.\ 1).

No created intellect, human or angelic, can reach the Divine Essence by its natural power (\textit{ST} I, q.\ 12, a.\ 4). Everything is known according to the mode of the knower. Our intellect naturally abstracts forms from sensible matter; an angel naturally knows immaterial substances whose existence is still participated. ``It follows therefore,'' Aquinas writes, ``that to know self-subsistent being is natural to the divine intellect alone; and this is beyond the natural power of any created intellect; for no creature is its own existence, forasmuch as its existence is participated. Therefore the created intellect cannot see the essence of God, unless God by His grace unites Himself to the created intellect, as an object made intelligible to it.''\endnote{St.\ Thomas Aquinas, \textit{Summa Theologiae} I, q.\ 12, a.\ 4.}

Nor can that vision ever occur through a created likeness impressed on the soul (\textit{ST} I, q.\ 12, a.\ 2). In ordinary knowing, a stone does not enter the mind in its physical substance; its intelligible likeness (\textit{species intelligibilis}) informs the intellect. With God, Aquinas gives three decisive reasons why no created likeness suffices: no lower order can represent a higher essence; God's essence \textit{is} His very act of existing (\textit{Ipsum Esse Subsistens}), whereas every created form differs from its own existence; and God is uncircumscribed, whereas every created form is finite. ``Hence to say that God is seen by some similitude,'' Aquinas concludes, ``is to say that the divine essence is not seen at all; which is false.''\endnote{St.\ Thomas Aquinas, \textit{Summa Theologiae} I, q.\ 12, a.\ 2; cf.\ \textit{Supplement}, q.\ 92, a.\ 1.}

In the \textit{Supplement} (q.\ 92, a.\ 1), Aquinas draws the sole conclusion left open: ``the Divine essence itself'' must be ``united to our intellect as its form, so as to be both that which is understood, and that whereby it is understood.'' God does not fuse with the creature into one natural substance; rather, strictly in the order of knowledge, His Essence actuates the intellect directly from within, without intermediary, so that ``the union of the body with the soul is an illustration of the blissful union of the spirit with God.''\endnote{Anonymous Compiler of the \textit{Supplementum Tertiae Partis} (drawn from St.\ Thomas Aquinas, \textit{Scriptum super Sententiis} IV, dist.\ 49), \textit{Supplement}, q.\ 92, a.\ 1, trans.\ Fathers of the English Dominican Province (1922).}

To bear that immediate union, God strengthens the finite mind by an infused supernatural disposition: the \textbf{Light of Glory} (\textit{lumen gloriae}), named from Psalm 35:10 (``In thy light we shall see light'') and Apoc 21:23 (``The glory of God hath enlightened it''):

\begin{quote}
Everything which is raised up to what exceeds its nature, must be prepared by some disposition above its nature; as, for example, if air is to receive the form of fire, it must be prepared by some disposition for such a form. But when any created intellect sees the essence of God, the essence of God itself becomes the intelligible form of the intellect. Hence it is necessary that some supernatural disposition should be added to the intellect in order that it may be raised up to such a great and sublime height. \dots\ By this light the blessed are made ``deiform''---i.e.\ like to God, according to the saying: ``When He shall appear we shall be like to Him, because we shall see Him as He is'' (1 John 3:2).\endnote{St.\ Thomas Aquinas, \textit{Summa Theologiae} I, q.\ 12, a.\ 5. The Gutenberg English Dominican printing prints ``1 John 2:2'' here, a typographical slip for 1 John 3:2.}
\end{quote}

Does this created light reintroduce a barrier between soul and God? Aquinas replies cleanly (\textit{ST} I, q.\ 12, a.\ 5, ad 2): the \textit{lumen gloriae} is not a medium \textit{in which} God is seen, like a mirror replacing His face, but a medium \textit{by which} the intellect is empowered to behold Him directly. The Church herself solemnly guarded the necessity of this supernatural elevation at the Council of Vienne (6 May 1312): in the constitution \textit{Ad nostrum}, the council condemned the Beghard error claiming that every intellectual nature is naturally blessed in itself and needs no light of glory raising it to behold and blessedly enjoy God.\endnote{Council of Vienne (6 May 1312), Constitution \textit{Ad nostrum} (DS 895).} Sanctifying grace is the seed; the light of glory is that same grace brought to its supernatural flower, rendering the creature \textit{deiform}.

Yet seeing God's Essence immediately does not exhaust the Infinite. The Church solemnly confesses as a dogma of faith that the one true God is \textit{incomprehensibilis} and \textit{ineffabilis} (Fourth Lateran Council, 1215, DS 800), infinite in intellect, will, and every perfection, and ineffably exalted above all things that exist or can be conceived outside Himself (\textit{super omnia quae praeter ipsum sunt et concipi possunt, ineffabiliter excelsus}; First Vatican Council, \textit{Dei Filius}, 24 April 1870, ch.\ 1, DS 3001).\endnote{Fourth Lateran Council (1215), Constitution 1 (\textit{Firmiter credimus}), DS 800; First Vatican Council (24 April 1870), Dogmatic Constitution \textit{Dei Filius}, ch.\ 1 (\textit{De Deo rerum omnium Creatore}), DS 3001.} In Question 12, Article 7, Aquinas distinguishes two senses of \textit{comprehension} (\textit{ST} I, q.\ 12, a.\ 7, ad 1; cf.\ I-II, q.\ 4, a.\ 3, ad 1). In the sense of attaining and holding the prize (1 Cor 9:24; Cant 3:4, ``I held him, and I will not let him go''), the blessed truly comprehend God. In the strict sense of \textit{inclusion}---knowing a thing as completely as it is knowable---no creature can ever comprehend God, because God is infinitely knowable while every created light of glory remains finite. The saints therefore see the whole God (\textit{totum}), for God is simple and holds back no hidden part, yet they never see Him exhaustively (\textit{totaliter}), in the infinite degree in which God knows Himself. That inexhaustibility seals heaven against satiety: the blessed dwell forever in an infinite Ocean of truth and love that fills their vessel while forever surpassing it.

What they behold is not a solitary Monad, but the living Trinity. Expounding John 17:3---``This is eternal life: That they may know thee, the only true God, and Jesus Christ, whom thou hast sent''---St.\ Augustine observes that the Holy Spirit is necessarily included, ``because He is the Spirit of the Father and Son, as the substantial and consubstantial love of both. For the Father and Son are not two Gods, nor are the Father and Son and Holy Spirit three Gods; but the Trinity itself is the one only true God.''\endnote{St.\ Augustine, \textit{Tractates on the Gospel of John} 105.3, trans.\ John Gibb and James Innes, NPNF1 vol.\ 7.} From that vision springs the inseparable delight of the will, as heat accompanies fire, for ``the very sight of God causes delight'' (\textit{ST} I-II, q.\ 4, aa.\ 1--2).

Both the \textit{Roman Catechism} and St.\ Bernard of Clairvaux reach for elemental imagery to express how this deifying vision transfigures the creature without destroying him. The \textit{Roman Catechism} writes that the blessed, ``whilst they retain their own nature, shall assume a certain admirable and almost divine form'':

\begin{quote}
As iron, when acted on by fire, becomes ignited, and, whilst it is substantially the same, seems changed into fire, which is a different substance; so the blessed, who are admitted into the glory of heaven, and who burn with a love of God, although they cease not to be the same, are yet affected in such a manner, as that they may be said with truth to differ more from the inhabitants of this earth, than iron, when ignited, differs from itself when cold.\endnote{\textit{Catechism of the Council of Trent}, Part I, Creed Article XII, trans.\ Donovan, pp.\ 95--96.}
\end{quote}

Four centuries earlier, in Chapter 10 of \textit{On the Love of God} (\textit{De diligendo Deo}), St.\ Bernard of Clairvaux described the fourth and highest degree of love---loving even oneself solely for God's sake---as the proper state of the blessed:

\begin{quote}
To attain to this, is for the soul to be deified; as a small drop of water appears lost if mixed with wine, taking its taste and colour; and as, when plunged into a furnace, a bar of iron seems to lose its nature and assume that of fire; or as the air filled with the sun's beams seems rather to become light than to be illuminated. So it is with the natural life of the Saints; they seem to melt and pass away into the will of God. For if anything merely human remained in man, how then should God be all in all? It is not that human nature will be destroyed, but that it will attain another beauty, a higher power and glory.\endnote{St.\ Bernard of Clairvaux, \textit{On the Love of God} (\textit{De diligendo Deo}), ch.\ 10, trans.\ Marianne Caroline and Coventry Patmore (London: Kegan Paul, Trench \& Co., 1884), p.\ 45.}
\end{quote}
